\patternSeven

\subsection*{Context}

\ac{bvi} students are increasingly represented in \ac{stem} (Science, Technology, Engineering, and Mathematics) degree programs. In computer science and engineering courses, learning is often driven by large collections of written materials such as lecture scripts, slide decks, exercise sheets, lab instructions, and reading assignments. \ac{bvi} students typically consume these materials through screen readers or Braille displays, which provide linear access; navigating and revisiting long documents is time-consuming.

Some students also benefit from alternative representations when sustained reading is challenging (e.g., dyslexia) or because initiating and sustaining attention on long texts is difficult (e.g., \ac{adhd}). For some autistic students, predictable structure and reduced cognitive overload can be beneficial, while others may prefer text over audio. These groups are secondary beneficiaries of this pattern; the primary accessibility motivation and constraints focus on \ac{bvi} learners.

This pattern applies both the teaching contexts with bundles of course materials and, with similar constraints, to academic work for a first-pass reading of long papers or book chapters.

\subsection*{Problem}

Long documents and heterogeneous material collections are costly to consume and revisit when access is primarily linear using screen readers or Braille displays. Students may lose time and context, enter topics late, and struggle to prepare for tutorials and discussions.

Students can use general-purpose AI systems to create summaries as easier accessible versions of documents. However, this ad-hoc approach requires repeated context building and prompt engineering for each new document and week. The accessibility workload is shifted to the student and does not scale.

In \ac{stem}, materials often include diagrams, plots, and schematics that are not inherently accessible. AI generated descriptions can support semantic orientation, but they do not replace precise accessible representations of visual artifacts.


\subsection*{Forces}
\renewcommand*\labelitemi{\textbullet}

\begin{itemize}
    \item Audio summaries are useful, but may contain omissions or inaccuracies. Quality assurance of the generated audios may become an additional effort.
    \item The audio should condense and restructure what is in the sources, but should not introduce novel aspects or even invented details.
    \item Some learners find audio distracting or unsuitable, hence alternatives must remain available.
    \item Pure audio is hard to skim. \ac{bvi} learners need additional information about the structure of documents like which chapter or timecodes. Ideally audio files are accompanied with  transcripts.
    \item Summaries must support learning without becoming a shortcut that replaces engagement with primary sources.
    \item Output quality depends on machine-readable, well-structured sources, but e.g. scan-only PDFs do not provide structural information.
    %\item Course materials and student submissions may be protected; institutional policies may restrict uploads or redistribution.
\end{itemize}

\begin{comment}
\begin{itemize}
    \item \textbf{Accuracy vs.\ effort:} Audio summaries are useful but may contain omissions or inaccuracies; quality assurance must not become a new large burden.
    \item \textbf{Faithfulness to sources:} The audio should condense and restructure what is in the sources, not introduce novel claims or invented details.
    \item \textbf{Audio helps many, not all:} Some learners find audio distracting or unsuitable; alternatives must remain available.
    \item \textbf{Navigability:} Pure audio is hard to skim; BVI learners need structure (chapters/timecodes) and ideally a transcript.
    \item \textbf{Assessment integrity:} Summaries must support learning without becoming a shortcut that replaces engagement with primary sources.
    \item \textbf{Source quality:} Output quality depends on machine-readable, well-structured sources (no scan-only PDFs; code should be text, not screenshots).
    \item \textbf{Privacy/IP constraints:} Course materials and student submissions may be protected; institutional policies may restrict uploads or redistribution.
\end{itemize}
\end{comment}
\subsection*{Solution}

\begin{sol}
Use a customized \ac{llm} tool specialized on generating overviews based on uploaded documents. 
Use the the audio overview feature of such tools to generate podcast-like, speech-first audio summaries from a curated bundle of course sources.
\end{sol}

Design the output to be speech-first, accessible, and faithful:
\begin{itemize}
    \item Start with learning objectives and a short overview.
    \item Define key terms and notation consistently.
    \item Provide  mini examples that illustrate the course style.
    \item Highlight common misconceptions.
    \item End with a short checklist (``You should be able to\ldots'') and ``what to read/do next.''
    \item Explicitly mark uncertainty and avoid adding details not supported by the provided sources.
\end{itemize}

Publish the audio as an  learning aid accompanied by additional information, especially a textual transcript and a chapter list with timecodes.
The audio is a learning aid and may contain errors or omissions, which should be declared in disclaimer.
%; primary sources remain authoritative.

Use different granularity levels to fit different needs, e.g., a short weekly overview, a longer deep dive for complex topics, and an exam-revision overview.

\subsection*{Resulting Context}

Instructors provide a reusable, low-friction first pass overview of course materials that reduces navigation overhead for \ac{bvi} students and supports preparation for tutorials and discussions. Students gain an additional channels, audio and transcript, that can be consumed flexibly  without replacing the original sources.

\subsection*{Examples}

The pattern was applied in the course ``Technical Informatics~1 (TI-1)'' at our university of applied sciences as a weekly pipeline: instructors curated a compact source bundle, generated an audio overview, and published audio, transcript and timecodes in the learning management system (LMS). In this case the used AI tool was NotebookLM\footnote{\url{https://notebooklm.google/}}. 

\begin{enumerate}
    \item Boolean operations and truth tables\\
    \emph{Input:} Text-first course content (e.g., Markdown-based core content rendered into slides) + and weekly exercise sheets.\\
    \emph{Request:} ``Create an Audio Overview with: (1) learning objectives, (2) key terms/notation, (3) one mini example, (4) two common misconceptions, (5) a checklist and next steps. Be faithful to the sources; mark uncertainty; avoid spatial references.''\\
    \emph{Output:} An audio overview, a transcript and chapter lists with timecodes covering AND/OR/NOT, truth-table construction, typical precedence pitfalls, and a short practice checklist.

    \item \ac{kv} maps \\
    \emph{Input:} Text-first KV-map representation generated from the exercise sheet.\\
    \emph{Request:} Same template, with emphasis on step-by-step strategy (grouping logic, minimization goals) and typical mistakes.\\
    \emph{Output:} Audio overview that provides intuition and a workflow for solving tasks. Audio output supports semantic orientation but does not replace a precise accessible KV-map representation.

    \item Circuit simulation interpretation\\
    \emph{Input:} Text-first simulation outputs tables of measured values with short additional explanations.\\
    \emph{Request:} Focus on interpretation: what metrics mean, how to compare cases, and how results link back to assumptions.\\
    \emph{Output:} Audio files and textual transcripts explaining which quantities matter, how to interpret changes, and a short checklist for reporting results.
\end{enumerate}

\subsection*{Consequences}

\noindent\textbf{Benefits}
\begin{itemize}
    \item[$+$] Podcast-Like audio summaries can provide a faster orientation over large material bundles. They reduce navigation overhead and support preparation for tutorials and discussions.
    \item[$+$]  Audio accompanied with textual information like a transcript and chapter list can be consumed flexibly and supports independent learning.
    \item[$+$] Students with dyslexia/ADHD benefit from a reduced reading load and lower entry barriers.
    %\item[$+$] \textbf{Scales for instructors.} A repeatable weekly pipeline yields consistent outputs with limited additional effort.
\end{itemize}

\noindent\textbf{Liabilities}
\begin{itemize}
    \item[$-$] AI generated materials may have inaccuracies and or e.g. misleading summaries. Quality assurance means an additional effort. Disclaimers should be added.
    \item[$-$] Important details of the original materials may be omitted due to condensation. Students have to consult primary sources for full fidelity.
    \item[$-$] Students may substitute audio summaries for primary sources, especially near assessments. Hence they may miss important details of the original course materials.
    \item[$-$] Audio can be distracting or unsuitable for some learners. Transcripts and text materials must remain available.
    \item[$-$]  Institutional policies may restrict what can be uploaded or redistributed.
\end{itemize}

\begin{comment}
\subsection*{Related Patterns}

\textbf{Pattern 6 (AI Tutor):} Use a course-specific, source-grounded tutor for precise transformations and explanations of visual artifacts (diagrams, schematics, KV maps). Pattern~7 complements this by providing orientation and revision via structured audio overviews.

\subsection*{Known Uses}
\begin{itemize}
    \item Applied in \textbf{Technical Informatics~1} as a weekly pipeline: instructors curate a short source bundle (script + accessible slide export + exercises), generate an Audio Overview, and publish audio + transcript + timecodes in the LMS.
    \item Used as a \textbf{first pass} before tutorials and as a \textbf{revision resource} during exam preparation.
\end{itemize}
\end{comment}

\subsection*{Related Work}

Structured educational podcasts have been studied as an accessible learning modality for \ac{bvi} users, showing that clear structure and navigability matter \cite{Buzzi2011StructuredPodcasts}. Speech-first summarization research emphasizes that summaries optimized for spoken output require different design choices than text summaries and should explicitly support linear comprehension \cite{McKeown2005TextToSpeech}. Recent accessibility research explores AI-assisted navigation and question answering to help \ac{bvi} learners access lecture content more efficiently \cite{Anderer2025LectureVideos}, aligning with our goal of reducing first-pass overhead for large material bundles. Recent studies of AI-generated educational podcasts report positive learner experiences and measurable learning outcomes in educational settings \cite{Do2024PAIGE}. 
Podcast-like course materials should be supplemented with textual information to enable targeted access to the audio information and to provide alternative channels.